\documentclass[11pt,a4paper]{letter}
\usepackage[utf8]{inputenc}
\usepackage[T1]{fontenc}
\usepackage[french]{babel}
\usepackage[left=2.5cm,right=2.5cm,top=2cm,bottom=2cm]{geometry}
\usepackage{hyperref}
\usepackage{xcolor}

% --- Couleurs CEA ---
\definecolor{primary}{RGB}{0, 82, 147}

\hypersetup{colorlinks=true, urlcolor=primary}

\signature{Loïc Aron MBASSI EWOLO}
\address{Loïc Aron MBASSI EWOLO\\
Cergy, France\\
(+33) 7 59 52 33 98\\
\href{mailto:loicaron.mbassiewolo@ensea.fr}{loicaron.mbassiewolo@ensea.fr}\\
\href{https://www.linkedin.com/in/loic-mbassi}{LinkedIn} | \href{https://github.com/Nameless0l}{GitHub}}

\begin{document}

\begin{letter}{CEA Grenoble\\
Laboratoire LSTA\\
17 avenue des Martyrs\\
38054 Grenoble Cedex 9}

\opening{Madame, Monsieur,}

Élève-ingénieur en 2A à l'ENSEA (double diplôme avec Polytechnique Yaoundé), je candidate au stage de 6 mois sur l'architecture GPGPU Vortex.

Côté FPGA, je travaille actuellement sur l'intégration d'un Nios V sur Cyclone V à l'ENSEA : synthèse Quartus, simulation ModelSim, et programmation C embarquée. Je suis à l'aise sous Linux et avec Python pour le scripting.

Au MILA cet été, j'ai développé des pipelines PyTorch pour analyser la convergence de modèles sur le papier de recherche d'\textbf{OpenAI} sur le \textbf{Grokking}. Le hackathon IndabaX (2e prix) m'a aussi appris à comparer rigoureusement différentes approches sur un même problème -- exactement ce que demande l'exploration de l'espace des configurations de Vortex.

Sur le challenge CoHoMa, j'ai travaillé sur Nvidia Jetson pour de l'IA embarquée. Ça m'a donné une première approche du calcul GPU, même si je n'ai pas encore manipulé d'architecture GPGPU au niveau RTL.

Je maintiens aussi un projet open-source (\textbf{Laravel API Generator}) qui génère automatiquement une architecture backend depuis des fichiers UML/XML. Ça m'a appris à structurer du code pour qu'il soit utilisable par d'autres.

Ce stage m'intéresse parce que je veux approfondir l'architecture des processeurs, et Vortex est un projet open-source concret pour ça. Je suis motivé pour apprendre ce qui me manque.

Disponible pour un entretien.

\closing{Je vous prie d'agréer, Madame, Monsieur, l'expression de mes salutations distinguées,}

\end{letter}
\end{document}
